\documentclass[10pt,a4paper]{amsart}
\usepackage[margin=19mm]{geometry}
\usepackage{amsmath,amssymb,mathtools}
\usepackage[scr=rsfs,cal=euler]{mathalfa}
\usepackage{xfrac}
\usepackage[unicode,hidelinks,bookmarks=false]{hyperref}
\newcommand{\ie}{i.e.\ }
\newcommand{\cf}{cf.\ }
\newcommand{\resp}{respectively\ }
\newcommand{\Subsection}{\subsection}
\providecommand{\nobreakdash}{\nobreak\hbox{-}}
\setlength{\emergencystretch}{2em}
\multlinegap0pt
\usepackage{tikz}
\usepackage{amssymb}
\usepackage[shortlabels]{enumitem}
\usepackage{mathtools}
\newtheorem{thm}{Theorem}
\newcommand{\pn}{pn}
\title{The number and average length of subpaths in graphs}
\author{Stijn Cambie}
\date{Source: arXiv:2608.23542v1 (CC BY 4.0)}
\begin{document}
\maketitle

\noindent\emph{Source and licence.} The number and average length of subpaths in graphs. Version arXiv:2608.23542v1 (CC BY 4.0), distributed by arXiv under the Creative Commons Attribution 4.0 International licence, http://creativecommons.org/licenses/by/4.0/, as recorded in the arXiv licence metadata for this exact version and shown on the abstract page https://arxiv.org/abs/2608.23542v1. Full licence notice: \texttt{licenses/arxiv-2608.23542-CC-BY-4.0.txt}.\par\smallskip

\noindent\textit{Editorial scope.} Counted unit: the article's thm thm (source0.tex lines 438--440). The article's complete original proof of it is delivered with it (source0.tex lines 442--526, one \texttt{proof} environment of 85 lines). No mathematics is written, repaired or re-derived: the statement and the proof are the article's own lines, in the article's own notation, and no retained line is substituted or re-flowed.  No further source-local context is needed: every cross-reference of every retained line resolves inside the two delivered spans.  Nothing was omitted from any retained span, so the package carries no omission record.  No retained line contains a citation, so the wrapper carries no bibliography and no bibliography entry.  No retained line contains a figure environment, an image-inclusion command or drawing code, so no image is delivered and the package declares no asset dependency.  Editorial formatting only, in addition: an AMS \texttt{amsart} preamble that restates the article's own declarations and macros used by the retained lines, namely the environments \texttt{thm} on separate counters, together with the standard packages those lines need.  The retained environments print in this document as Theorem 1 in order of appearance, whereas the article prints the same items with the numbering of its own full document.  The complete native source of the article as distributed in the arXiv e-print for this exact version is bundled unchanged as \texttt{sources/208417/source0.tex}.
\begin{thm}
    Among triangle-free graphs, the balanced complete bipartite graph maximizes (uniquely) the number of subpaths.
\end{thm}

\begin{proof}
Write $T_n=K_{\lfloor n/2\rfloor,\lceil n/2\rceil}$ (it is a Turan graph)
and let $m_s(G)$ denote the number of matchings of size $s$ in $G$.

We first record that, among triangle-free graphs of order $n$ and for every $s\geq0$,
\begin{equation}\label{eq:matching-bound}
        m_s(G)\le m_s(T_n).
\end{equation}

We prove~\eqref{eq:matching-bound} by induction on $s$. The case $s=0$ is trivial. Every
matching of size $s$ is counted once for each of its $s$ edges, and
therefore
\[
 s\,m_s(G)
   =\sum_{uv\in E(G)}m_{s-1}(G-\{u,v\}).
\]
Since $G-\{u,v\}$ is triangle-free, induction and Mantel's theorem give
\[
 s \cdot m_s(G)
 \leq e(G)m_{s-1}(T_{n-2})
 \leq
 \left\lfloor\frac{n^2}{4}\right\rfloor
 m_{s-1}(T_{n-2})= s \cdot m_s(T_n).
\]


We now bound the paths using matchings. First consider a path of odd
length $2s-1$. Orient the path and take the edges in positions
\[
        1,3,\ldots,2s-1.
\]
These edges form a matching of size $s$, together with an ordering of
its edges.

Fix an ordered matching $e_1,\ldots,e_s$. There are at most two ways
to orient its edges so that they form a path in the prescribed order.
Indeed, there are two choices for the orientation of $e_1$. Once the
terminal vertex of $e_i$ has been chosen, it is adjacent to at most
one endpoint of $e_{i+1}$, since adjacency to both endpoints would
create a triangle with $e_{i+1}$.

Consequently, the number of oriented paths of length $2s-1$ is at most $2s!\,m_s(G).$
Every unoriented path has two orientations, so
\begin{equation}\label{eq:odd-path}
        \pn_{2s-1}(G)\leq s!\,m_s(G).
\end{equation}

For paths of even length $2s$, use the same matching formed by the
first $2s$ vertices and then append the last vertex. For a fixed
ordered matching, there are again at most two compatible orientations.
If both exist, their terminal vertices are the two endpoints $x,y$ of
the last matching edge. Since $xy\in E(G)$ and $G$ is triangle-free,
\[
        N(x)\cap N(y)=\emptyset.
\]
Hence the total number of unused vertices that can extend these two
orientations is at most $n-2s$. The same bound is immediate if only
one orientation exists. Therefore
\begin{equation}\label{eq:even-path}
        \pn_{2s}(G)
        \leq \frac{s!}{2}(n-2s)m_s(G).
\end{equation}

Both \eqref{eq:odd-path} and \eqref{eq:even-path} are equalities for
every complete bipartite graph: for each ordered matching there are
exactly two compatible orientations, and in the even case the numbers
of possible final vertices add up to $n-2s$. Combining these
equalities with \eqref{eq:matching-bound}, we obtain
\[
        \pn_\ell(G)\leq\pn_\ell(T_n)
\]
for every $0\leq\ell\leq n-1$. Summing over $\ell$ gives
\[
        \pn(G)\leq\pn(T_n).
\]

Finally, equality in the sum forces equality for paths of length one.
Thus
\[
        e(G)=\pn_1(G)=\pn_1(T_n)
             =\left\lfloor\frac{n^2}{4}\right\rfloor.
\]
The equality case of Mantel's theorem implies $G\cong T_n$. Hence the
balanced complete bipartite graph is the unique maximizer.
\end{proof}

\end{document}
